


\section{Rappels conception relationnelle}

\begin{frame}
\frametitle{Rappel sur la conception relationnelle}

Pour l'essentiel.

\begin{redbullet}
  \item déterminer les "entités" (film, réalisateurs, acteurs) pertinentes pour l'application; 
  \item définir une méthode d'identification de chaque entité;  
  \item préserver le lien entre les entités.
\end{redbullet}

\vfill

Exemple pour notre application "Les films"

\begin{redbullet}
  \item  On distingue  deux types d'entités: les films et les réalisateurs.
  \item On les identifie par un numéro incrémental (l'id)
  \item On référence, dans chaque film, le réalisateur grâce à son identifiant.
\end{redbullet}


\end{frame}



\begin{frame}[fragile]
\frametitle{Illustration}

\alert{La table des films}

\medskip

\begin{small}
   \begin{tabular}[!ht]{l|l|l}
    \textbf{id} & \textbf{titre} & \textbf{année} \\
     \hline
    1 & Alien & 1979 \\
    2 & Vertigo & 1958 \\
    3 & Psychose & 1960 \\
    4 & Kagemusha & 1980 \\
    5 & Volte-face & 1997 \\
    6 & Pulp Fiction & 1995 \\
     \hline
   \end{tabular}
   \end{small}
 \medskip
 
\alert{La table des artistes}
\medskip

\begin{small}
   \begin{tabular}[!ht]{l|l|l|l}
    \textbf{id} & \textbf{nom} & \textbf{prénom} &  \textbf{année} \\
     \hline
    101 & Scott & Ridley & 1943 \\
    102 & Hitchcock & Alfred & 1899 \\
    103 & Kurosawa & Akira & 1910 \\
    104 & Woo & John & 1946 \\
    105 & Tarantino & Quentin & 1963 \\
     \hline
   \end{tabular}
 \end{small}
\medskip

\alert{Il manque le lien entre les films et les artistes.}

\end{frame}
 
 
\begin{frame}[fragile]
\frametitle{Le système de référencement en relationnel}

\alert{La table des films.}

\medskip

 \begin{small}
   \begin{tabular}[!ht]{l|l|l|l}
    \textbf{id} & \textbf{titre} & \textbf{année} & \textbf{idRéalisateur} \\
     \hline
    1 & Alien & 1979 & \alert{101} \\
    2 & Vertigo & 1958 & \alert{102}\\
    3 & Psychose & 1960 & \alert{102} \\
    4 & Kagemusha & 1980 & \alert{103} \\
    5 & Volte-face & 1997 & \alert{104} \\
    6 & Pulp Fiction & 1995 & \alert{105} \\
     \hline
   \end{tabular}
   \end{small}
 \medskip
 
\alert{La table des artistes}
\medskip

\begin{small}
   \begin{tabular}[!ht]{l|l|l|l}
    \textbf{id} & \textbf{nom} & \textbf{prénom} &  \textbf{année} \\
     \hline
    \alert{101} & Scott & Ridley & 1943 \\
    \alert{102} & Hitchcock & Alfred & 1899 \\
    \alert{103} & Kurosawa & Akira & 1910 \\
    \alert{104} & Woo & John & 1946 \\
    \alert{105} & Tarantino & Quentin & 1963 \\
     \hline
   \end{tabular}
\end{small}
\medskip

\alert{Ce n'est pas un lien physique, mais il permet des jointures.}

\end{frame}
 
 
\begin{frame}
\frametitle{Méthode générale: schémas entité/association}

Exemple pour notre base de films.

\vfill

\figSlideWithSize{films}{9cm}

\end{frame}


\begin{frame}[fragile]
\frametitle{Schémas relationnels}

Exemple pour notre base de films.

\vfill
\begin{minted}{sql}
create table Artiste  (idArtiste integer not null,
                       nom varchar (30) not null,
                       prenom varchar (30) not null,
                       anneeNaiss integer,
                       primary key (idArtiste),
                       unique (nom, prenom))
\end{minted}

\vfill

Exprime des \alert{contraintes}: clé primaire, unicité (nom, prénom), typage, valeurs obligatoires.

\end{frame}


\begin{frame}[fragile]
\frametitle{Schémas relationnels (2)}

La table des films:

\vfill
\begin{minted}{sql}
create table Film  (idFilm integer not null, 
                   titre    varchar (50) not null,
                   annee    integer not null,
                   idRealisateur    integer not null,
                   genre varchar (20) not null,
                   resume      varchar(255),
                   codePays    varchar (4),
                   primary key (idFilm),
                   foreign key (idRealisateur) 
                       references Artiste);
\end{minted}

\vfill
La \alert{contrainte d'intégrité référentielle} \texttt{foreign key} garantit la cohérence
du référencement.

\end{frame}


\begin{frame}[fragile]
\frametitle{La normalisation relationnelle}

En relationnel, toutes les données sont "à plat". Cela tend à multiplier le nombre
de tables contenant les informations sur une même entité.

\vfill

Exemple: pour représenter toutes les informations sur un film, il faut

\begin{redbullet}
  \item La table \texttt{Film} (une ligne pour le film)
  \item La table \texttt{Artiste} (une ligne pour le réalisateur, autant de lignes que d'acteurs)
  \item La table \texttt{Role} (autant de lignes que de rôles)
  \item et la table \texttt{Pays}, et d'autres encore...
\end{redbullet}

\vfill

Il faut beaucoup \alert{d'écritures}, et autant de \alert{lectures} pour reconstituer
l'information.

\end{frame}


\begin{frame}{Représentation relationnelle du film (Pulp Fiction)}

\begin{columns}
    \begin{column}{0.2\textwidth}
Les films:
\end{column}
  \begin{column}{0.7\textwidth}
   \begin{tabular}[!ht]{l|l|l|l}
      id & titre & année & idReal\\
     \hline
    17 & Pulp Fiction & 1994 & 37\\
     \hline
   \end{tabular}
   \end{column}
   \end{columns}
   
\bigskip

\begin{columns}
    \begin{column}{0.2\textwidth}
Les artistes: 
\end{column}
  \begin{column}{0.7\textwidth}

   \begin{tabular}[!ht]{l|l|l|l}
    id & nom & prénom & naissance\\
     \hline
    37 & Tarantino & Quentin & 1963\\
    11 & Travolta & John & 1954\\
    27 & Willis & Bruce & 1955\\
     \hline
   \end{tabular}
\end{column}
\end{columns}

  \bigskip

\begin{columns}
    \begin{column}{0.2\textwidth}
Les rôles
\end{column}
  \begin{column}{0.7\textwidth}

   \begin{tabular}[!ht]{l|l|l}
    idFilm & idArtiste & rôle\\
     \hline
    17 & 11 & Vincent Vega \\
    17 & 27 & Butch Coolidge\\
   17 & 37 & Jimmy Dimmick\\
     \hline
   \end{tabular}
  \end{column}
\end{columns}

\end{frame} % ending frame Exemple de structure relationnelle


\begin{frame}[fragile]{Impact de la normalisation}

\begin{remblock}{Essentiel: effet de la normalisation}
\begin{redbullet}
  \item il faut effectuer \alert{plusieurs écritures} pour une même entité (transactions)
  \item il faut effectuer des \alert{jointures} pour reconstituer une entité.
\end{redbullet}
\end{remblock}

\vfill

Reconstitution de l'information sur un film en SQL:

\begin{minted}{sql}
select titre, nom, prenom, role
from Film, Artiste, Role
where role='Vincent Vega'
and Film.id = Role.idFilm
and Artiste.id = Role.idActeur
\end{minted}

\end{frame}


%%%%%%%%%%%%
\section{Modélisation NoSQL}


\begin{frame}[fragile]
\frametitle{Puissance du modèle structuré}

Les documents structurés ne sont pas soumis aux contraintes de normalisation.


\vfill

Un attribut peut avoir \alert{plusieurs valeurs} (en utilisant la structure de tableau)
\begin{minted}[baselinestretch=.9,
               fontsize=\small,
               framesep=2mm]{javascript}
    {
      "title": "Pulp fiction",
      "year": "1994",
      "genre": ["Action", "Policier", "Comedie"]
      "country": "USA"
    }
\end{minted}


\vfill

En relationnel, il faudrait ajouter deux tables.


\end{frame}

\begin{frame}[fragile]
\frametitle{Puissance du modèle structuré (2)}

Il est également facile de représenter des données régulières. 

\medskip

   \begin{tabular}[!ht]{l|l|l}
    \textbf{id} & \textbf{nom} & \textbf{prénom} \\
     \hline
    37 & Tarantino & Quentin \\
     \hline
    167 & De Niro & Robert \\
     \hline
    168 & Grier & Pam \\
     \hline
   \end{tabular}
   
 \vfill
 
Facile à représenter sous forme d'un document (très régulier).

\medskip
   
\begin{minted}[baselinestretch=.9,
               fontsize=\small,
               framesep=2mm]{javascript}
 [
   artiste:  {"id": 37, "nom": "Tarantino", "prenom": "Quentin"},
   artiste:  {"id": 167, "nom": "De Niro", "prenom": "Robert"},
   artiste:  {"id": 168, "nom": "Grier", "prenom": "Pam"} 
 ] 
\end{minted}

\vfill

Table relationnelle = arbre de hauteur constante. Trois niveaux: ligne, attribut, valeur.

\end{frame}

\begin{frame}[fragile]
\frametitle{Plus puissant que la représentation relationnelle?}


\begin{remblock}{On peut donc aussi représenter des données régulières}
Probablement pas du tout efficace. \alert{Pourquoi\,?}
\end{remblock}

$\Rightarrow$ peut être utile pour échanger des informations\,; pas pour les
stocker.

\vfill

\begin{beamerboxesrounded}{Premier constat}

\begin{redbullet}
  \item Dans une base relationnelle, les données sont très contraintes / structurées: \alert{on peut stocker
  séparément la structure (le schéma) et le contenu (la base)}.
  \item Une représentation arborescente XML / JSON est plus appropriée pour des données de structure
       \alert{complexe} et / ou \alert{flexible}.
\end{redbullet}
\end{beamerboxesrounded}


\end{frame}


\begin{frame}[fragile]
\frametitle{Documents structurés = imbrication des structures}

Grâce à \alert{l'imbrication des structures}, il est possible de
représenter dans une même unité un film \alert{et} son metteur en scène.

\smallskip

\begin{minted}[baselinestretch=.9,
               fontsize=\small,
               framesep=2mm]{javascript}
{ "title": "Pulp fiction",
    "year": "1994",
    "genre": "Action",
    "country": "USA",
    "director": {"last_name": "Tarantino",
	               "first_name": "Quentin",
	               "birth_date": "1963"	
               }
}
\end{minted}

\smallskip

Dans les bases documentaires\,: on essaie de créer des unités d'information \alert{autonomes} 
pour limiter les écritures et éviter d'avoir à faire des jointures.

\begin{remblock}{Intérêt?}
Les systèmes NoSQL sont conçus pour passer à l'échelle par distribution. C'est en
partie incompatible avec les jointures et les transactions.
\end{remblock}
\end{frame}



\begin{frame}[fragile]
\frametitle{Représentation sous forme de document structuré}

On peut coder toutes les informations sur un film.
\begin{minted}[baselinestretch=.9,
               fontsize=\footnotesize,
               framesep=2mm]{javascript}
{   "_id": "movie:17",
    "title": "Pulp Fiction",
    "year": "1994",
    "director": 	{
	    "last_name": "Tarantino", "first_name": "Quentin", 
	    "birth_date": "1963"	
	    },
    "actors": [
       {
         "first_name": "John", "last_name": "Travolta", 
         "birth_date": "1954",  "role": "Vindent Vega"	
       },
       {
        "first_name": "Bruce", "last_name": "Willis", 
        "birth_date": "1955", "role": "Butch Coolidge"	
       },
       {
        "first_name": "Quentin", "last_name": "Tarantino", 
        "birth_date": "1963", "role": "Jimmy Dimmick"	
       }
     ]
}
\end{minted}

\end{frame}


\begin{frame}[fragile]{Bilan - à retenir}

Les avantages de la représentation par document structuré.

\begin{redbullet}

  \item \alert{Plus besoin de jointure (?)} : il est inutile de faire des jointures pour reconstituer
    l'information puisqu'elle n'est plus dispersée, comme en relationnel, dans plusieurs
    tables.
   
   \item \alert{Plus besoin de transaction (?)}: une écriture (du document) suffit; 
    une lecture suffit pour récupérer l'ensemble des informations.
    
  \item \alert{Adaptation à la distribution}. Si les documents sont autonomes, il est très facile
    des les déplacer pour les répartir au mieux dans un système distribué.
\end{redbullet}

\vfill

Un système NoSQL est beaucoup plus facile à implanter qu'un système relationnel. Une
opération \textit{put()}, une opération \textit{get()} et c'est parti.

\end{frame}


\begin{frame}[fragile]{Les inconvénients du modèle}

La représentation par document a deux inconvénients forts. 

\begin{redbullet}
  \item \alert{Chemin d'accès privilégié}: les   films apparaissent près de la racine des documents, les artistes 
    sont enfouis dans les profondeurs;
    
     L'accès aux films est donc privilégié

 \item \alert{Les entités ne sont plus autonomes}: 
   pas moyen de créer un réalisateur s'il n'y a pas au moins un film.

    
  \item \alert{Redondance}: la même information doit être représentée
    plusieurs fois, ce qui est tout à fait fâcheux (Quentin Tarantino est représenté deux fois).
 \end{redbullet}
 
\vfill

La redondance mène à des incohérences.

\vfill

Privilégier un chemin d'accès est bon pour certaines applications, mauvais pour d'autres.

\vfill

Comment faire pour obtenir la liste des films de Quentin Tarantino avec
la représentation précédente?
\end{frame}

\begin{frame}[fragile]{Les autres inconvénients}

\alert{Pas de langage de requête?}

Il faut faire un programme pour tout, même la moindre mise à jour!

\vfill
\alert{Pas de schéma}?

On peut mettre n'importe quoi dans la base; c'est l'application qui doit faire les contrôles et les corrections.

\vfill

\alert{Pas de transaction?}

Ne convient pas pour beaucoup d'applications (commerce électronique).


\end{frame}

\section{Bilan, NoSQL ou relationnel?}

\begin{frame}[fragile]
\frametitle{Bilan, NoSQL ou relationnel?}

Quand utiliser (ou pas) une base documentaire\,?

\begin{redbullet}
  \item \alert{Des données très spécifiques, peu ou faiblement structurées} (texte, données multimédia, graphes)
  
  \item \alert{Peu de mises à jour, beaucoup de lectures.}; la redondance ne pose pas de problème.
  
  \item \alert{on veut traiter de très gros volumes de manière ``scalable''}.     


  \item \alert{De forts besoins en temps réel}.     
\end{redbullet}

Conditions non remplies\,: un système relationnel est probablement une bien meilleure option.

\begin{remblock}{Bien retenir}
NoSQL = \textit{Not Only SQL}. Personne (de sérieux) ne prétend que ces systèmes vont
remplacer les systèmes relationnels, sauf pour certaines niches.
\end{remblock}

\end{frame}

\begin{frame}[fragile]
\frametitle{Quelques éléments de réflexion}


 \alert{Toujours se poser (au moins) les questions suivantes}:

\begin{redbullet}
  \item S'agit-il de représenter une structure régulière (p.e. une table)?
  
    Dans ce cas il n'est pas nécessaire d'intégrer la structure et le contenu.
        
  \item Est-ce que je dois souvent modifier les données existantes?
  
    Dans ce cas j'ai sans doute besoin de transactions, et je risque de payer cher toute redondance.
    
  \item La structure d'un document est-elle fixe et prévisible?
  
  Si oui, le relationnel est bien adapté.
\end{redbullet}

\vfill

\begin{remblock}{La principale niche des bases NoSQL}
Comme entrepôts de données où on accumule des informations dans l'optique de traitements analytiques.
\end{remblock}


\end{frame}



